\begin{center}
\large{\textbf{LEMBAR PENGESAHAN}} \\
\vspace{4mm}
\textbf{PERFORMA MESIN PANAS KUANTUM FRAKSIONAL BERBASIS PARTIKEL PADA KOTAK 1 DIMENSI} \\
\vspace{4mm}
%\large{\textbf{SCALAR FIELD POTENTIAL MODEL OF THE INFLATIONARY COSMOLOGY}} \\
%\vspace{5mm}
\textbf{TUGAS AKHIR} \\
\vspace{4mm}
Diajukan Untuk Memenuhi Salah Satu Syarat \\
Memperoleh Gelar Sarjana Sains \\
pada Bidang Fisika Teori \\
Program Studi S-1 Departemen Fisika \\
Fakultas Sains dan Analitika Data\\
Institut Teknologi Sepuluh Nopember Surabaya \\
\vspace{4mm}
Oleh : \textbf{PUTRA NAUFAL TABASSAM} \\
\vspace{1mm}
NRP. 5001221120 \\
\vspace{4mm}
\begin{center}
	Disetujui oleh Tim Penguji Tugas Akhir:
\end{center}

\vspace{4ex}

\begingroup
% Menghilangkan padding
\setlength{\tabcolsep}{0pt}

\noindent
\begin{tabularx}{\textwidth}{X l}
	% Ubah kalimat-kalimat berikut dengan nama dosen pembimbing pertama
	1. Dr. Heru Sukamto, M.Si.     & Pembimbing 1 \\
	NIP 198411092012121001       &  \\
	&  \\
	&.........................  \\
	% Ubah kalimat-kalimat berikut dengan nama dosen pembimbing kedua
	2. Dr. Gagus K. Sunnardianto   & Pembimbing 2\\
	NIP: 198611052019021001 	&
	................................... \\
	&  \\
	&  \\
	% Ubah kalimat-kalimat berikut dengan nama dosen penguji pertama
	3. Malik Anjelh Baqiya, M.Si., Ph.D.   & Penguji I \\
	NIP 198210202008121003     & \\
        &  \\
	& ......................... \\
        &  \\
	&  \\
	% Ubah kalimat-kalimat berikut dengan nama dosen penguji kedua
	4. Prof. Drs. Agus Purwanto, M.Si, M.Sc, D.Sc.    & Penguji II \\
	NIP 196408111990021001     &  \\
	&  \\
        & .........................\\
	% Ubah kalimat-kalimat berikut dengan nama dosen penguji ketiga
\end{tabularx}
\endgroup
  \begin{center}
	\textbf{SURABAYA\\Juli, 2024}
\end{center}
\end{center}

\newpage
\thispagestyle{empty}
\vspace*{\fill}
    \begin{center}
	Halaman ini sengaja dikosongkan
    </center}
\vspace*{\fill}
\pagebreak
